\chapter{إشارة الإنذار}
% full version: chapters/14_pain.tex

\pencilfig{14}{\pencilart{14}}{الألم جهاز إنذار بالحريق، لا مقياس حرارة.}

كل الحواس الأخرى تخبر كيف هو العالم. أما الألم فلا. إنه يخبر أن شيئًا يهدد الجسم، وعلوّ صوته تحدده الآلة نفسها إلى حد كبير، لا المنبه. الألم جهاز إنذار بالحريق، لا مقياس حرارة.

مستشعرات الألم لا تعمل إلا على القوي: على حرارة تتجاوز الأربعين درجة وبضع درجات، وعلى ضغط شديد، وعلى مواد كاوية. وهي تعتاد أقل بكثير من مستشعرات اللمس: جهاز الإنذار يجب ألا يصمت ما دام التهديد قائمًا.

\section*{سلكان}

تذكّر حين صدمت إصبعك. أولًا ألم حاد ودقيق: أين ومتى. وبعد لحظة ألم كليل منتشر طويل: ما مدى السوء. هذان سلكان مختلفان. أحدهما سريع، تصل إشارته في عشرات الأجزاء من ألف من الثانية. والآخر بطيء، تسير إشارته نحو ثانية.

\section*{البوابة}

لماذا نفرك موضع الصدمة بعد أن نصطدم؟ في الحبل الشوكي تمر إشارة الألم عبر ”بوابة“. واللمس يوقظ عصبونًا كابحًا يوارب البوابة. لذلك يساعد فرك الكدمة فعلًا، لكن ضد الألم المعتدل فقط؛ أما الشديد فيمر على أي حال. وأن الألم الشديد يعطّل بنفسه العصبون الكابح عندئذ جزء من نظرية البوابة الأصلية، ولا تأكيد مباشرًا له حتى الآن.

\pencilfig{126}{%
% gate: the pain wire drives the relay neuron; touch wakes an inhibitory neuron that closes the gate
\draw[pen=pcAmber] (0,1.35) -- (3.05,1.0); \draw[arr=pcAmber] (2.8,1.03) -- (3.08,1.0);
\node[lbl, text=pcAmber!70!black, anchor=east] at (-0.05,1.35) {ألم};
\draw[pen=pcBlue] (0,-0.15) -- (1.75,-0.15); \draw[arr=pcBlue] (1.5,-0.15) -- (1.78,-0.15);
\node[lbl, text=pcBlue, anchor=east] at (-0.05,-0.15) {لمس};
\draw[dotpen=pcRed, wash=pcRed] (2.05,-0.15) circle (0.26);
\draw[pcRed, line width=0.7pt, -{Bar[width=6pt]}] (2.25,0.05) -- (3.2,0.66);
\node[lbl, text=pcRed, anchor=north] at (2.05,-0.45) {كابح};
\draw[dotpen=pcGray, wash=pcGray] (3.4,0.9) circle (0.34);
\node[lbl, anchor=south] at (3.4,1.3) {البوابة};
\draw[pen] (3.75,0.9) -- (5.6,0.9); \draw[arr] (5.35,0.9) -- (5.65,0.9);
\node[lbl, anchor=west] at (5.7,0.9) {إلى الدماغ};
}{البوابة في الحبل الشوكي. تمر إشارة الألم إلى الدماغ عبر عصبون البوابة، واللمس يوقظ عصبونًا كابحًا يواربها. لذلك يساعد فرك الكدمة، لكن ضد الألم المعتدل فقط.}

\section*{مقبض الصوت من الأعلى}

يستطيع الدماغ بنفسه أن يرفع صوت الألم أو يخفضه. من الأعلى، من الدماغ، تنزل إلى البوابة قنوات إذاعية: \nt{SERO} و\nt{NORA} ومواد خاصة تشبه المسكنات يفرزها الدماغ بنفسه. في موقف خطر قد يكاد الألم يختفي: من يهرب من الخطر لا وقت لديه ليعرج. وترقّب الراحة يخفف الألم أيضًا، جزئيًا عبر المواد نفسها.

\section*{حين يعلق جهاز الإنذار}

الألم المتكرر يزيد الحساسية: الصدمة نفسها تبدأ بإيلام أكثر. وهذا أيضًا تعلم: حلقة تشد نفسها بنفسها ما دام الجرح يلتئم. لكن الحلقة القوية فيها خطر: قد ”تعلق“ وتواصل الرنين بعد أن يلتئم الجرح، كمجموعة العصبونات من الفصل الثاني التي تغذي نفسها بنفسها.

\section*{المعلم والمقاطعة}

والألم معلم أيضًا: بالنسبة إلى ”أفضل أو أسوأ مما كان متوقعًا“ الدوبامينية يعمل مكافأة سالبة. وهو يقاطع العمل الجاري، كرشقة النورأدرينالين: اترك كل شيء، واهتم بالخطر. وأسرع رد فعل، سحب اليد عن الساخن، يُغلق مباشرة في الحبل الشوكي بزوج العضلات المتضادة نفسه من الفصل السابق، قبل أن تلحق أن تشعر بالألم.

\fullversion{14}
